\vspace{-0.5em}
\section{Conclusion}\label{sec:conclusion}
This paper presented \emph{PaTAS}, a Subjective Logic framework that propagates trust in neural networks through a parallel structure of \emph{Trust Nodes} and \emph{Trust Functions} running alongside feedforward and backpropagation. Its \emph{Parameter Trust Update} and \emph{Inference-Path Trust Assessment} mechanisms capture how input quality, learned parameters, and activation paths jointly shape the trustworthiness of a prediction.

Experimental evaluation on real-world and adversarial datasets shows that PaTAS detects trust degradation caused by data poisoning and adversarial patching while maintaining interpretability and stability, providing reliability information not reflected by accuracy alone, in particular when the provenance of the training data is unknown. In quantitative comparison with maximum-softmax, MC~Dropout, and Evidential Deep Learning, PaTAS is the only method to reliably detect out-of-distribution and corrupted inputs (AUROC $0.975$ and $0.97$ vs.\ near-chance baselines) and to flag training-label corruption without any test data, while matching the softmax filter on in-distribution selective prediction after identical calibration; its applicability to a given dataset is predictable a priori from the conformity model's own statistics (\cref{tab:diagnostics}).

Several directions remain open, which we state explicitly. The evaluation used shallow fully-connected architectures on small benchmarks; extending trust propagation to deeper convolutional and transformer models, and enriching the input-trust construction beyond marginal conformity so that unregistered imagery (GTSRB, CIFAR-10) moves inside the applicability boundary, are the primary next steps. Three further studies would strengthen the framework: a sensitivity analysis of how PaTAS degrades when the initial input and label opinions are themselves noisy or incorrect; a deeper study of the gradient-to-evidence mapping beyond the $\epsilon$ sensitivity reported in \cref{tab:cancer_results}; and an empirical measurement of runtime and memory overhead against the analytical $3\times$ estimate. The security scope could also broaden from detection to provable robustness, with guarantees on trust degradation under bounded poisoning and evaluation against evasion, model-stealing, and membership-inference attacks. On the modeling side, the effect of activation functions on trust propagation is uncharacterized, and the gradient-counting used for \(T_{\theta_i \mid y_\text{batch}}\) fixes the uncertainty component by construction, motivating adaptive quantification schemes that also account for parameter magnitude.
